\chapter{Inputs and outputs}\label{ch:c6}
This chapter is the reference for every command and every output, followed by a complete example set-up.

\begin{plain}
This is the look-up chapter: every command you can write in the input files, what it does and its default, and every output
file with its columns. There is no need to read it from start to finish; come here when you need the exact spelling of a
command or the meaning of an output column.
\end{plain}

\section{Input reference}\label{ch:inputs}
This chapter lists every command. Commands are case sensitive; separators may be commas or spaces, as in all Raven files.

\subsection{\texttt{.rvi} (main input file)}
\begin{table}[!htbp]\centering
\caption{Commands in the \cmd{.rvi} file.}
\begin{tabular}{P{0.34\linewidth}P{0.608\linewidth}}\toprule\hdr
Command & Meaning\\\midrule
\cmd{:GroundwaterModel MODFLOW6} & switches the coupling on (\cmd{NONE}: off, the default)\\
\cmd{:rvg\_Filename file} & name of the groundwater file, relative to the \cmd{.rvi} (default \cmd{<model>.rvg})\\
processes into \cmd{GROUNDWATER} & any process moving water into \cmd{GROUNDWATER} in a coupled HRU produces recharge; use
\cmd{:-->Conditional HRU\_GROUP} so uncoupled HRUs keep their own baseflow\\
\cmd{:Transport} & not available with \cmd{:GroundwaterModel MODFLOW6}: the coupling exchanges water only, no
constituent mass (recharge would leave its mass in Raven's \cmd{GROUNDWATER} store, and seepage and groundwater discharge
would return water without mass), so a coupled run with any \cmd{:Transport} command (including \cmd{AGE} and
\cmd{TEMPERATURE}) stops with a message naming the constituent\\\bottomrule
\end{tabular}
\end{table}

\subsection{\texttt{.rvh} (HRUs and subbasins)}
The existing \cmd{AQUIFER\_PROFILE} column of \cmd{:HRUs} holds a profile name, or \cmd{[NONE]}. Reservoirs to be coupled
need \cmd{:SeepageParameters} and a lake HRU (\cmd{:HRUID}) with an \cmd{AQUIFER\_PROFILE}.

\subsection{\texttt{.rvp} (parameters)}
\begin{lstlisting}
:AquiferClasses
  :Attributes, K_HORIZ, K_VERT, SPEC_STORAGE, SPEC_YIELD, POROSITY
  :Units,      m/d,     m/d,    1/m,          none,       none
  GRAVEL,      50.0,    5.0,    1.0E-5,       0.25,       0.30
  CLAY_TILL,   0.001,   0.0001, 1.0E-4,       0.03,       0.40
:EndAquiferClasses
:AquiferProfiles
  # name, number of layers, then {class, thickness [m] or TO_BEDROCK, layer type} for each layer
  VALLEY_ALLUVIUM, 3, GRAVEL, 15.0, AQUIFER, CLAY_TILL, 5.0, AQUITARD, SAND, TO_BEDROCK, CONFINED_AQUIFER
  UPLAND_TILL,     1, UPLAND_TILL, TO_BEDROCK, AQUIFER
:EndAquiferProfiles
:AquiferProfileParameters
  :Attributes,     INITIAL_HEAD_DEPTH, SEEPAGE_LEAKANCE, RECHARGE_DELAY
  :Units,          m,                  1/d,              d
  VALLEY_ALLUVIUM, 3.0,                1.0,              0.0
:EndAquiferProfileParameters
\end{lstlisting}
\begin{table}[!htbp]\centering
\caption{Aquifer class attributes (\cmd{.rvp}).}
\begin{tabular}{P{0.27\linewidth}P{0.13\linewidth}P{0.528\linewidth}}\toprule\hdr
Aquifer class attribute & Default & Meaning\\\midrule
\cmd{K\_HORIZ} & required & horizontal hydraulic conductivity [m/d]\\
\cmd{K\_VERT} & $K_h/10$ & vertical hydraulic conductivity [m/d]\\
\cmd{SPEC\_STORAGE} & $10^{-5}$ & specific storage [1/m]\\
\cmd{SPEC\_YIELD} & 0.1 & specific yield [--]\\
\cmd{POROSITY} & 0.3 & saturated water content $\theta_s$ of the unsaturated columns [--] (\cmd{:SoilZoneVadose}; at least \cmd{SPEC\_YIELD})\\
\cmd{CLAPP\_B} & 4 & Clapp--Hornberger exponent $b$ of the unsaturated columns [--]\\
\cmd{AIR\_ENTRY\_PRESSURE} & 0.3 m & air-entry suction $\psi_a$ of the unsaturated columns [m]\\\bottomrule
\end{tabular}
\end{table}
Layer types: \cmd{AQUIFER} (or \cmd{UNCONFINED\_AQUIFER}), \cmd{CONFINED\_AQUIFER}, \cmd{AQUITARD},
\cmd{AQUICLUDE} (not as top layer). Classes must be defined before the profiles that use them.
With \cmd{:MF6Simulation} a profile may have no layers (\cmd{name, 0}); the model's own layers are used.
\begin{table}[!htbp]\centering
\caption{Aquifer profile parameters (\cmd{.rvp}).}
\begin{tabular}{P{0.3\linewidth}P{0.13\linewidth}P{0.498\linewidth}}\toprule\hdr
Profile parameter & Default & Meaning\\\midrule
\cmd{INITIAL\_HEAD\_DEPTH} & 5 m & initial depth to water below land surface\\
\cmd{DEFAULT\_BEDROCK\_DEPTH} & 50 m & bedrock depth where no bedrock map is given\\
\cmd{SOIL\_ZONE\_DEPTH} & 0 m & model top below land surface\\
\cmd{SEEPAGE\_LEAKANCE} & 1 d$^{-1}$ & seepage conductance per m$^2$ of cell\\
\cmd{SEEPAGE\_SMOOTHING\_DEPTH} & 0.5 m & depth below the top where seepage starts\\
\cmd{RIVERBED\_K} & 0.5 m/d & riverbed hydraulic conductivity\\
\cmd{RIVERBED\_THICKNESS} & 1 m & riverbed thickness\\
\cmd{EXTINCTION\_DEPTH} & 0 m & water-table ET extinction depth (0 = off)\\
\cmd{RECHARGE\_DELAY} & 0 d & recharge delay time constant (0 = off)\\\bottomrule
\end{tabular}
\end{table}

\subsection{\texttt{.rvg} (groundwater file)}
\begin{longtable}{P{0.45\linewidth}P{0.498\linewidth}}
\caption{Commands in the \cmd{.rvg} file.}\\
\toprule\hdr
Command & Meaning (default)\\\midrule\endfirsthead
\caption[]{Commands in the \cmd{.rvg} file (continued).}\\
\toprule\hdr
Command & Meaning (default)\\\midrule\endhead
\cmd{:HRUGeometry file \{ID field\}} & HRU polygons, GeoJSON in lon/lat (or \cmd{:InputCRS}); ID field (\cmd{HRU\_ID}). \textbf{Required.}\\
\cmd{:RiverGeometry file \{subbasin field\}} & river lines; without a field, subbasins are assigned automatically\\
\cmd{:LandSurface file} & DEM, ESRI ASCII grid or GeoTIFF (\cmd{.tif}, \cmd{.tiff}) in lon/lat or \cmd{:InputCRS}; Section~\ref{sec:rasters} (HRU \cmd{ELEVATION})\\
\cmd{:InputCRS EPSG:code | TMERC | ALBERS | LCC \dots} & projection of every geometry and raster input file; wells may then be given by \cmd{X}, \cmd{Y} (lon/lat); Section~\ref{sec:inputcrs}\\
\cmd{:BedrockSurface file} & bedrock elevation, same formats (\cmd{DEFAULT\_BEDROCK\_DEPTH})\\
\cmd{:GridCellSize m} & cell size (1000 m)\\
\cmd{:MinCellCoverage f} & fraction of a cell that coupled HRUs must cover (0.25)\\
\cmd{:GridType REGULAR|QUADTREE|HRU\_MESH|FILE} & kind of grid (\cmd{REGULAR}); Section~\ref{sec:gridtypes}\\
\cmd{:MF6Simulation file} & couple an existing MODFLOW 6 simulation (Chapter~\ref{ch:existing})\\
\cmd{:MF6Model}, \cmd{:MF6StartDate} & model to couple (if several); date of the model's time zero\\
\cmd{:MF6CRS}, \cmd{:MF6GridFile}, \cmd{:OverlapWeights} & how HRUs are linked to the model's cells (one of them)\\
\cmd{:RavenTakesOver}, \cmd{:KeepPackages} & recharge/ET/UZF packages replaced by Raven or kept\\
\cmd{:StreamPackage p \{AUX a|DOMINANT\_HRU\}}, \cmd{:SeepageToRaven ON|OFF} & stream package to Raven's reaches; seepage package\\
\cmd{:GridRotation} $\theta$ & rotation of regular and quadtree grids, degrees counterclockwise (0)\\
\cmd{:GridRefinement RIVERS|WELLS|LAKES s}, \cmd{LINES|POLYGONS file s} & finer cells (size $s$ [m]) along rivers, around wells, in and around reservoir lakes, or along/around features in a file\\
\cmd{:GridFile file \{ID\}} & imported cell polygons, with \cmd{:GridType FILE} (ID field \cmd{CELL\_ID})\\
\cmd{:QuadtreeMaxLevel n} & maximum halvings of the base cell (from the refinement sizes)\\
\cmd{:FlowCorrection AUTO|XT3D|XT3D\_RHS|NONE} & XT3D conductance tensor (\cmd{AUTO}: off; \cmd{XT3D\_RHS} or \cmd{XT3D} to switch it on)\\
\cmd{:LayerConnection INDEX|ELEVATION} & connect layers by number or by overlapping elevation (\cmd{INDEX})\\
\cmd{:MF6Library path} & MODFLOW 6 library (else the environment variable \cmd{RAVEN\_MF6\_LIB}, next to the Raven executable, then \cmd{lib/mf6/} of the Raven checkout, then the system path; Section~\ref{sec:getmf6})\\
\cmd{:SeepageReturnTo SV} & Raven store receiving seepage (lowest soil layer; soil-less HRUs: \cmd{SURFACE\_WATER})\\
\cmd{:RiverBedDepth m | REFERENCE\_FLOW} & riverbed depth below the DEM (\cmd{REFERENCE\_FLOW})\\
\cmd{:GWInitialization STEADY\_STATE r} & steady-state start with recharge $r$ mm/d (off)\\
\cmd{:SurfaceWater\allowbreak Exchange SV group LEAKANCE L} & Raven store of an HRU group $\leftrightarrow$ aquifer (also with \cmd{:MF6Simulation})\\
\cmd{:ReservoirExchange} & Raven reservoirs $\leftrightarrow$ aquifer (also with \cmd{:MF6Simulation})\\
\cmd{:Wells} \dots\ \cmd{:EndWells} & table: ID, \cmd{NAME, LATITUDE, LONGITUDE, SCREEN\_TOP, SCREEN\_BOT} (with \cmd{:InputCRS}: \cmd{X, Y} instead of \cmd{LATITUDE, LONGITUDE})\\
\cmd{:ObservationWells} \dots & table: ID, \cmd{NAME, LATITUDE, LONGITUDE, SCREEN\_ELEV} (or \cmd{X, Y})\\
\cmd{:GeneralHeadBoundary name file HEAD h CONDUCTANCE c \{LAYERS ALL|TOP\}} & head-dependent boundary; $c$ per cell [m$^2$/d]\\
\cmd{:SpecifiedHeadBoundary name file HEAD h \{LAYERS ALL|TOP\}} & fixed-head boundary\\
\cmd{:SolverMax\allowbreak OuterIterations n} & Newton iterations per step (500)\\
\cmd{:SolverHead\allowbreak Tolerance m} & Newton head-change tolerance (0.00001 m)\\
\cmd{:HeadSaveFrequency n} & steps between saves of the full head field (30)\\
\cmd{:WriteNetCDFHeads} & write \cmd{GWHeads.nc} (Raven compiled with NetCDF)\\
\cmd{:WriteGWBudgetFile} & MODFLOW budget file \cmd{mf6/gwf.cbc} and \cmd{mf6/gwf.porosity} for particle tracking (off; generated models; Section~\ref{sec:budgetfile})\\
\cmd{:LinkageCache file} & linkage cache location (\cmd{<output>/mf6/linkage.cache})\\
\cmd{:SoilZone CONCEPTUAL | PHYSICAL SOIL[0] \dots} & the soil zone and the water table (\cmd{CONCEPTUAL}); with \cmd{PHYSICAL} the further \cmd{:SoilZone*} commands of Table~\ref{tab:szcmds} apply\\\bottomrule
\end{longtable}

\subsection{Soil-zone commands (\texttt{.rvg})}\label{sec:szcmdref}
These commands apply only with \cmd{:SoilZone PHYSICAL} (Section~\ref{ch:soilzone}); without it they are reported as
ignored. A command given where it has no effect (\cmd{:SoilZoneVadoseClasses} or \cmd{:SoilZoneVadoseInitialFlux} with
\cmd{:SoilZoneVadose NONE}, \cmd{:SoilZoneAquiferClass} on a model Raven builds) is reported as ignored too. Values
outside the ranges given stop the run with a message.
\begin{longtable}{P{0.36\linewidth}P{0.15\linewidth}P{0.42\linewidth}}
\caption{Soil-zone commands in the \cmd{.rvg} file.}\label{tab:szcmds}\\
\toprule\hdr Command and syntax & Default & Meaning\\\midrule\endfirsthead
\caption[]{Soil-zone commands in the \cmd{.rvg} file (continued).}\\
\toprule\hdr Command and syntax & Default & Meaning\\\midrule\endhead
\cmd{:SoilZone CONCEPTUAL} \newline \cmd{:SoilZone PHYSICAL SOIL[0] \dots\ SOIL[$M-1$]} & \cmd{CONCEPTUAL} & \cmd{PHYSICAL}: the listed soil layers (from the top, consecutive) are the soil zone, owned by Raven; MODFLOW starts below them and the water table may rise into them (Section~\ref{ch:soilzone}). The last \cmd{:SoilZone} counts\\
\cmd{:SoilZoneIteration n \{$\epsilon_h$\} \{tol$_Q$\}} & 100, 0.0001~m, 0.000001 & interface updates after MODFLOW's outer iterations, $n\ge0$ (0: explicit coupling); soil-level tolerance [m] and relative flow tolerance, $>0$ (Section~\ref{sec:szinterface})\\
\cmd{:SoilZoneSubgrid n} & 5 & DEM points per cell side ($n^2$ points per cell) for the saturated fraction on sloping and irregular land, $1\le n\le20$ (Section~\ref{sec:szslope})\\
\cmd{:SoilZoneOutput group} & all soil-zone HRUs & HRU group written to \cmd{GWSoilExchange.csv}\\
\cmd{:SoilZoneVadose NONE} \newline \cmd{:SoilZoneVadose RICHARDS \{$\epsilon$\} \{N\}} \newline \cmd{:SoilZoneVadose KINEMATIC \{$\epsilon$\} \{N\}} \newline \cmd{:SoilZoneVadose LAG \{$t_m$\} \{N\}} & \cmd{NONE}; $\epsilon$ 0.02, $t_m$ 0, $N$ 16 & unsaturated columns between the soil-zone base and a deeper water table: Richards' equation, kinematic wave or linear storage cascade (Sections~\ref{sec:szvadose}, \ref{sec:szvadosekw}, \ref{sec:szvadoselag}); $\epsilon$ largest content change in a substep, $(0,0.2]$; $t_m$ travel time per metre [d/m], $\ge0$ (0: derived from each material); $N$ nodes per column at most, 3 to 200. \cmd{NONE}: recharge reaches the water table in the step\\
\cmd{:SoilZoneVadoseClasses K} & 1 & unsaturated columns per exchange column, one per depth class of its DEM points, $1\le K\le10$ (Section~\ref{sec:vzclasses})\\
\cmd{:SoilZoneVadoseInitialFlux q} & 0 & on a cold start, every unsaturated column starts at its steady profile for the flux it receives, $(1-F)\,q$, $q\ge0$ [mm/d]; 0: at rest (reference profile); ignored on a groundwater hotstart\\
\cmd{:SoilZoneCapillaryRise ON | OFF} & \cmd{ON} with \cmd{RICHARDS}, \cmd{KINEMATIC} or \cmd{LAG}; \cmd{OFF} with \cmd{NONE} & capillary rise from the water table (and, with \cmd{RICHARDS}, from the columns' top nodes) into soil-zone layers below field capacity (Section~\ref{sec:szcaprise})\\
\cmd{:SoilZoneExchangeTiming LAGGED | AFTER} & \cmd{LAGGED} & \cmd{LAGGED}: last step's exchange is lent to the soil zone before Raven's processes and settled after the step's exchange; \cmd{AFTER}: the exchange follows Raven's processes (Section~\ref{sec:sztiming})\\
\cmd{:SoilZoneClasses K} & 1 & Raven's processes run on up to $K$ wetness classes of each HRU's soil-zone columns, $1\le K\le50$; needs \cmd{:Method ORDERED\_SERIES} (Section~\ref{sec:szclasses})\\
\cmd{:SoilZoneAquiferClass class} & $b$ 4, $\psi_a$ 0.3~m, $\theta_s$ $S_y+0.1$ & existing models only: the aquifer class whose \cmd{CLAPP\_B}, \cmd{AIR\_ENTRY\_PRESSURE} and \cmd{POROSITY} the unsaturated columns use (conductivity and specific yield come from the model; a steady model without specific yield takes the class's \cmd{SPEC\_YIELD} when given, else 0.2; Section~\ref{sec:szexisting})\\\bottomrule
\end{longtable}
The soil classes of the soil-zone layers need \cmd{POROSITY}, \cmd{FIELD\_CAPACITY} $<1$, and, for the bottom layer,
\cmd{HYDRAUL\_COND} [mm/d]; capillary rise uses their \cmd{CLAPP\_B} and \cmd{AIR\_ENTRY\_PRESSURE} [mm]. Recharge must leave the
soil zone from \cmd{SOIL[$M-1$]} into \cmd{GROUNDWATER}, and \cmd{RECHARGE\_DELAY} must be 0 for soil-zone HRUs (a delay through a
deep unsaturated zone: \cmd{:SoilZoneVadose LAG}, Section~\ref{sec:szvadoselag}). The
aquifer profile parameter \cmd{SEEPAGE\_SMOOTHING\_DEPTH} is the contact depth $d_s$ of the soil zone (below 0.1~m a
warning); \cmd{SOIL\_ZONE\_DEPTH} and \cmd{SEEPAGE\_LEAKANCE} are not used for soil-zone HRUs.
The block commands take one row per item, for example:
\begin{lstlisting}
:Wells
  :Attributes, NAME, LATITUDE, LONGITUDE, SCREEN_TOP, SCREEN_BOT
  1, PW_1, 59.27310, -124.73460, 1199.7, 1188.7
:EndWells
:ObservationWells
  :Attributes, NAME, LATITUDE, LONGITUDE, SCREEN_ELEV
  101, OW_NEAR, 59.27310, -124.73060, 1192.7
:EndObservationWells
:GeneralHeadBoundary   GHB1  boundary_line.geojson HEAD 1310.0 CONDUCTANCE 5000.0 LAYERS TOP
:SpecifiedHeadBoundary LAKE1 lake.geojson          HEAD 1312.0 LAYERS TOP
:SurfaceWaterExchange  DEPRESSION WetlandHRUs LEAKANCE 0.01
:ReservoirExchange
\end{lstlisting}

\subsection{\texttt{.rvt} (time series)}
\begin{lstlisting}
:WellFlowRate 1 m3/d                  # well ID; negative = pumping
  1985-10-01 00:00:00 1.0 1826
  0.0
  ...
:EndWellFlowRate
:BoundaryHead GHB1 m              # replaces HEAD of boundary GHB1
  1985-10-01 00:00:00 1.0 1826
  ...
:EndBoundaryHead
:ObservationData GW_HEAD 101 m    # observation-well ID; head AT each date
  1985-10-01 00:00:00 1.0 1826
  ...
:EndObservationData
\end{lstlisting}

\subsection{\texttt{.rvc} (initial conditions)}
\par\noindent\begin{minipage}{\linewidth}
\begin{lstlisting}
:InitialGWHeads DEPTH_BELOW_SURFACE 5.0     # or FROM_PROFILE (default) or ELEVATION z
\end{lstlisting}
\end{minipage}\par\smallskip
\cmd{:GWHeads nlay nrow ncol}, \cmd{:GWRechargeStore}, \cmd{:GWRiverLossCarry} and \cmd{:GWReservoirSeepage}, and with
\cmd{:SoilZone PHYSICAL} \cmd{:GWSoilZoneColumns}, \cmd{:GWVadoseColumns}, \cmd{:GWVadoseClassColumns} and \cmd{:GWSoilZoneLag},
are written by Raven into \cmd{solution.rvc} for hotstart; they are read back automatically (Section~\ref{app:hotstart}).

\section{Output reference}\label{ch:outputs}
All rows are stamped with the date at the \emph{end} of the step (the moment the values refer to).
\subsection{\texttt{GWBudget.csv} --- one row per step}
\begin{table}[!htbp]\centering
\caption{Columns of \cmd{GWBudget.csv}.}
\begin{tabular}{P{0.43\linewidth}P{0.518\linewidth}}\toprule\hdr
Column & Meaning\\\midrule
time, date & model time [d] and date at the end of the step\\
Raven recharge [m$^3$/d] & recharge leaving Raven (before any delay)\\
MF6 recharge [m$^3$/d] & recharge entering MODFLOW\\
seepage to Raven [m$^3$/d] & drain discharge at the model top\\
storage release [m$^3$/d] & water released from aquifer storage (negative: stored)\\
MF6 balance error [m$^3$/d], [\%] & $\varepsilon$ of Section~\ref{ch:balance}, and relative to gross flow\\
returned to Raven [m$^3$] & seepage volume added to Raven stores\\
outer iterations, converged & Newton iterations used; 1 target tolerance reached, 2 accepted at the relaxed tolerance, 0 not converged\\
recharge delay storage [m$^3$] & water in delay reservoirs\\
recharge onto CHD cells [m$^3$/d] & recharge of HRUs lying entirely on fixed-head cells\\
river leakage to aquifer, aquifer discharge to rivers [m$^3$/d] & river exchange, each direction (on a step MODFLOW did not converge, without \cmd{:SoilZone PHYSICAL}: the discharge handed to Raven, after the withholding of Section~\ref{sec:guard}; so are \emph{seepage to Raven} and \emph{aquifer discharge to surface stores})\\
river exchange applied to Raven [m$^3$] & net volume actually added to (or taken from) Raven reaches\\
river loss carried to next step [m$^3$] & river loss not supplied in the step, taken from the river the next step (normally 0)\\
water-table ET [m$^3$/d] & evaporation from the water table\\
wells specified, wells actual [m$^3$/d] & requested and delivered well rates\\
GHB, CHD [m$^3$/d] & boundary flows (positive into the aquifer)\\
surface-store leakage / discharge [m$^3$/d], applied [m$^3$] & wetland and lake exchange\\
reservoir seepage sent, delivered [m$^3$/d], awaiting delivery [m$^3$] & reservoir exchange\\
reservoir gain not supplied by MODFLOW [m$^3$] & a gain the lake received that MODFLOW did not supply (a drying cell, or cells under the lake that were specified heads or inactive on the day it was delivered); taken back from the water reaching the subbasin in the same step (normally 0; a warning and a run total)\\
other model packages [m$^3$/d] & with an existing model: net flow of the model's own packages (wells, boundaries, \dots), positive into the aquifer\\
soil zone to aquifer, aquifer to soil zone [m$^3$/d] & \cmd{:SoilZone PHYSICAL}: soil-zone exchange, each direction\\
rejected recharge, soil-zone overflow, exchange applied to Raven, shortfall [m$^3$] & Section~\ref{ch:soilzone}; shortfall is 0 (except on steps MODFLOW ended unconverged at a low outer-iteration limit)\\
interface updates, converged, change [m] & soil-zone interface iteration of the step\\
further soil-zone columns & layer redistribution, withheld discharge, negative recharge not held, unsaturated zone (inflow, to aquifer, capillary rise, storage), capillary rise from the water table, lagged-exchange debt and repayment, specified-head exchange, recharge not applied, reservoir seepage not delivered: Section~\ref{app:outputs}\\\bottomrule
\end{tabular}
\end{table}

\noindent One row of \cmd{GWBudget.csv} from the 20-year Liard run (selected columns):
\begin{center}\small
\begin{tabular}{ll}\toprule\hdr
Column & Value\\\midrule
\cmd{time [d]}, \cmd{date} & 3910, 1996-06-15\\
\cmd{Raven recharge [m3/d]} & 21\,365\,037.27\\
\cmd{MF6 recharge [m3/d]} & 21\,365\,037.27\\
\cmd{seepage to Raven [m3/d]} & 19\,087\,056.98\\
\cmd{storage release [m3/d]} & $-3\,925\,200.44$\\\bottomrule
\end{tabular}
\end{center}
Recharge leaving Raven equals recharge entering MODFLOW; on this wet June day the aquifer stores about 3.9 million m$^3$.

\subsection{Other files}
\begin{table}[!htbp]\centering
\caption{Other output files.}\label{tab:outputs}
\begin{tabular}{P{0.25\linewidth}P{0.698\linewidth}}\toprule\hdr
File & Content\\\midrule
\cmd{GWHRUState.csv} & per coupled HRU: overlap-weighted mean water table of its columns and mean depth to water\\
\cmd{GWSoilExchange.csv} & \cmd{:SoilZone PHYSICAL}: per soil-zone HRU and step, soil-zone base, soil water level and water table (area means over all the HRU's soil-zone columns, those over specified-head cells included, whose water table is the specified head; empty for an HRU without a column), saturated fraction, exchange with the aquifer both ways (\emph{soil zone to aquifer}, \emph{aquifer to soil zone}: with active cells, as \cmd{GWBudget.csv} books it), rejected recharge, overflow, soil water and capacity; in a model with specified heads also \emph{soil zone to specified heads} and \emph{specified heads to soil zone} (columns over specified-head cells and links whose cell is a specified head in the step, with the capillary rise a specified head supplies). Summed over the HRUs, each exchange column equals the column of the same name in \cmd{GWBudget.csv} (Section~\ref{ch:soilzone}). Two further columns are model-wide values of the step, repeated on every HRU's row: \emph{interface updates} and \emph{interface change [m]} (as in \cmd{GWBudget.csv})\\
\cmd{GWHeads.csv} & head at each monitoring well, from $t=0$ (missing value $-1.2345$ when the screened cell is dry)\\
\cmd{GWConnectivity.csv} & number of separate wet areas; wet fraction of each profile\\
\cmd{mf6/grid\_cells.geojson} & every cell as a lon/lat polygon: number, active, profile, dominant HRU, area\\
\cmd{GWModelSummary.txt} & MODFLOW 6 library and its version, grid type and size, MODFLOW format, HRU fidelity of the cells, active and pass-through cells, profiles, river/well/boundary/store/reservoir cells, linkage per HRU (areas, fraction linked, cells), whether the linkage came from the cache\\
\cmd{GWHeads.nc} & (optional) head(time, layer, y, x) with 2-D latitude and longitude, CF conventions\\
\cmd{mf6/} & the generated MODFLOW model, its list file (\cmd{gwf.lst}, \cmd{mfsim.lst}) and binary heads (\cmd{gwf.hds}); with \cmd{:WriteGWBudgetFile} also the budget file \cmd{gwf.cbc} and \cmd{gwf.porosity}\\
\cmd{Diagnostics.csv} & Raven's usual file, with \cmd{GW\_HEAD} statistics\\
\cmd{Watershed\allowbreak Storage.csv} & Raven's balance; includes all exchanged water\\
\cmd{solution.rvc} & includes \cmd{:GWHeads}, \cmd{:GWRechargeStore}, \cmd{:GWReservoirSeepage}, \cmd{:GWRiverLossCarry}, \cmd{:GWSoilZoneColumns}, \cmd{:GWVadoseColumns}, \cmd{:GWVadoseClassColumns}, \cmd{:GWSoilZoneLag}\\\bottomrule
\end{tabular}
\end{table}

\subsection{Budget file and particle tracking}\label{sec:budgetfile}
With \cmd{:WriteGWBudgetFile} a generated model saves MODFLOW's cell-by-cell flows: the model's name file gets
\cmd{SAVE\_FLOWS} (every package), the NPF package \cmd{SAVE\_SPECIFIC\_DISCHARGE} and \cmd{SAVE\_SATURATION}, and the output
control writes \cmd{mf6/gwf.cbc} at the steps it saves the heads (\cmd{:HeadSaveFrequency}). \cmd{mf6/gwf.porosity} holds the
\cmd{POROSITY} of each cell's aquifer class in MODFLOW's array layout, written as the \cmd{GRIDDATA} block of a PRT \cmd{MIP6} file.
Nothing else changes: without the option the \cmd{mf6} folder is the same, and with it the solution is identical.

These files let MODFLOW~6's particle-tracking model (PRT) or MODPATH run on the coupled flows after the Raven run.
\cmd{tools/make\_prt.py OUTPUT\_FOLDER [-{}-backward] [-{}-points file] [-{}-run libmf6]} builds a PRT simulation in
\cmd{OUTPUT\_FOLDER/prt/}: release points around each well of \cmd{:Wells} (or from a file), the flow model's periods, and
the head and budget files, reversed with FloPy for backward tracking (capture zones and times of travel). With \cmd{-{}-run}
it runs PRT through the MODFLOW~6 library and writes \cmd{endpoints.csv} (travel time and end point of every particle). PRT
tracks particles on DIS and DISV grids (regular and quadtree grids); for tracking through time, save the flows every step
(\cmd{:HeadSaveFrequency 1}). An existing model keeps its own output control: add \cmd{BUDGET FILEOUT} and
\cmd{SAVE BUDGET} to its OC file instead (Chapter~\ref{ch:existing}).

\section{A complete example set-up}
The following \cmd{.rvg} uses most options (the soil-zone commands of Table~\ref{tab:szcmds} keep their defaults except those shown); lines not needed can be deleted. With \cmd{:SoilZone PHYSICAL} the \cmd{:SurfaceWaterExchange} group must not name a soil-zone layer (its \cmd{DEPRESSION} store lies on the ground above the soil zone, and its exchange has its bed there, Section~\ref{ch:stores}), and the \cmd{.rvi} takes recharge from \cmd{SOIL[1]}.
\begin{lstlisting}
# ---- geometry ----
:HRUGeometry        my_hrus.geojson   HRU_ID
:RiverGeometry      my_rivers.geojson SubId
:LandSurface        dem_lonlat.asc
:BedrockSurface     bedrock_lonlat.asc
# ---- grid and solver ----
:GridCellSize       1000.0
:MinCellCoverage    0.25
:SolverHeadTolerance 0.00001
:SolverMaxOuterIterations 500
:HeadSaveFrequency  30
:WriteNetCDFHeads
:LinkageCache       cache/linkage.cache
# ---- coupling options ----
:MF6Library         /opt/mf6/libmf6.so
:SeepageReturnTo    SOIL[2]
:RiverBedDepth      REFERENCE_FLOW
:GWInitialization   STEADY_STATE 0.5
:SurfaceWaterExchange DEPRESSION WetlandHRUs LEAKANCE 0.01
:ReservoirExchange
# ---- soil zone and water table (optional) ----
:SoilZone PHYSICAL SOIL[0] SOIL[1]
:SoilZoneVadose RICHARDS
:SoilZoneVadoseClasses 3
:SoilZoneClasses 3
# ---- stresses and observations ----
:Wells
  :Attributes, NAME, LATITUDE, LONGITUDE, SCREEN_TOP, SCREEN_BOT
  1, TownWell, 49.123, -80.456, 245.0, 230.0
:EndWells
:ObservationWells
  :Attributes, NAME, LATITUDE, LONGITUDE, SCREEN_ELEV
  101, OW_1, 49.120, -80.450, 240.0
:EndObservationWells
:GeneralHeadBoundary   East  east_edge.geojson HEAD 210.0 CONDUCTANCE 2000.0 LAYERS ALL
:SpecifiedHeadBoundary Lake  lake.geojson      HEAD 255.0 LAYERS TOP
\end{lstlisting}
with, in the \cmd{.rvt},
\begin{lstlisting}
:WellFlowRate 1 m3/d
  2000-01-01 00:00:00 1.0 3
  -1500.0
  -1500.0
  -1800.0
:EndWellFlowRate
:ObservationData GW_HEAD 101 m
  2000-01-01 00:00:00 1.0 3
  241.20
  241.18
  -1.2345
:EndObservationData
\end{lstlisting}
(\cmd{-1.2345} marks a missing observation, as elsewhere in Raven). Only three days are shown: like every Raven
forcing series, a \cmd{:WellFlowRate} series must cover the whole simulation, and its time interval must divide evenly
into one day (Raven stops with a message otherwise). Observation series may be shorter.
